En un problema de transporte clásico, supongamos que el coste de envío entre un origen $i$ y un destino $j$ incluye un coste fijo por activar la ruta (por ejemplo, contratar un camión) más un coste variable por unidad de producto transportado. Además, existe una restricción que impide enviar suministros a más de $k$ destinos en total desde cualquier origen. 

¿Por qué el método del Simplex de transporte ya no es válido para resolver este problema?

\textbf{Solución}

En el caso de que haya un coste fijo por envío, sería necesario definir una variable binaria. En este caso, no se podría resolver el problema con el método del simplex y no podríamos disponer porque se trataría de un problema lineal con varaibles enteras (y, en particular, binarias).
